\begin{graphicspathcontext}{{./chapters/ai/imgs/},{./chapters/ai/imgs/auto/},\old}

\sidenote{Image générée avec Mistral AI}
\begin{frame}[t]{{Systèmes multiagents :} Concepts clés}
	\vspace{-.2cm}
	\begin{columns}
		\begin{column}{.8\linewidth}
			\begin{block}{Qu'est-ce qu'un agent ? \cite{Wooldridge09}}
				Une \Emph{entité autonome humaine, matérielle ou logicielle} qui :
				\begin{compactdescription}
				\item[Perçoit] son environnement {\smaller(capteurs, messages, flux de données)}
				\item[Raisonne] à l'aide de connaissances internes et d'objectifs
				\item[Agit] pour poursuivre ses objectifs {\smaller(actionneurs, appels de services, décisions)}
				\item[Communique] et \Emph{coopère} avec d'autres agents
				\end{compactdescription}
			\end{block}
			\begin{block}{Qu'est-ce qu'un système multiagents ? \cite{Wooldridge09}}
				\begin{compactdescription}
				\item[Décentralisé] aucun point de contrôle ou de défaillance unique
				\item[Évolutif] ajout/suppression d'agents dynamique
				\item[Hétérogène] agents ont des rôles, compétences et connaissances différents
				\item[Comportement émergent] résultats complexes issus de règles locales simples
				\end{compactdescription}
			\end{block}
		\end{column}
		\begin{column}{.2\linewidth}
			\includegraphics[width=\linewidth]{mas_application} \\[1cm]
			\begin{center}
				
\includegraphics[width=.8\linewidth]{sarllogo} \\
				{\smaller Framework SARL}\\
				\url{www.sarl.io}
			\end{center}
		\end{column}
	\end{columns}
\end{frame}

\begin{frame}[t]{{IA agentique :} Un système multiagents simplifié et spécialisé}
	\alertbox{IA agentique $=$ raisonnement LLM \textbf{+} principes de coordination des systèmes multiagents}
	\small
	\begin{columns}
		\begin{column}{.5\linewidth}
			\begin{block}{Orchestration centrée sur les LLM}
				\begin{compactdescription}
				\item[LangChain / LangGraph] chaînes d'appels LLM avec
				utilisation d'outils, mémoire et graphes d'agents cycliques
				\item[OpenAI Agents SDK] SDK léger pour construire des
				pipelines d'agents avec transfert et appels d'outils
				\item[AutoGen (Microsoft)] framework multi-agents conversationnel ;
				les agents négocient et s'autocorrigent
				\item[CrewAI] équipes d'agents basées sur des rôles avec
				délégation de tâches et collaboration explicites
				\end{compactdescription}
			\end{block}
    		\end{column}
		\begin{column}{.5\linewidth}
			\begin{stabularx}[title={Choisir le bon outil}]{X|X}
			\tabularhead{Besoin}{Frameworks possibles} \\
			Prototypage rapide & LangChain, LangGraph, CrewAI \\
			\hline
			Systèmes robustes & \textit{SARL}, \textit{JADE} \\
			\hline
			Simulation et analytique & \textit{Gamma}, \textit{Repast}, \textit{SARL} \\
			\hline
			IoT / edge distribué & \textit{SPADE}  \\
			\end{stabularx}
			{\smaller\smaller \textit{Italique} : framework multiagents généraliste}
    		\end{column}
	\end{columns}
	\normalsize
\end{frame}

\end{graphicspathcontext}

\endinput

